\scanpage{27}
\uline{Pb} Les composantes connexes de $\prod_{\unsure{X/k}}P/X$ sont elles de type fini sur $k$ ?

\uline{Corollaire 2} Supposons, \deleted{\unsure{pour}} \inserted{pour} \deleted{\unsure{les définitions}} $P/X/k$, $X$ propre sur $k$, $P/X$ loc. de type fini, $k$ corps (pas \unsure{néc.} \inserted{\unsure{sép.}} alg. clos) que \deleted{$\forall$} pour $\forall\ x_i\in\operatorname{Ass}\mathcal O_X$, $\exists\ x'_i\in\bar x$ tel que $k(x'_i)/k$ radiciel. \inserted{[plus gén., tel que $\prod_{k(x'_i)/k}P_{x'_i}/k(x'_i)$ représentable]} Alors $\prod_{X/k}P/X$ représentable.
\margin{\deleted{\unsure{différent du} \unsure{précédent ?} \unsure{car les} \illegible{2} \unsure{sont les mêmes} : \unsure{l'un} \illegible{2} \unsure{pt de} $\unsure{P}$ \illegible{1}}}
\margin{Si les $P_{x'_i}$ sont quasi-projectifs, alors $\prod_{X/k}P/X$ est essentiellement quasi-projectif.}
P. ex. si $X$ \deleted{est} est \deleted{\unsure{un schéma} irréductible} \inserted{\unsure{irréductible|irrédundant}}, et admet un pt $x$ radiciel sur $k$ \deleted{Et}, $\prod_{X/k}P/X$ existe.

\uline{Corollaire 3} $P/X/k$ comme dessus. Alors $\exists\ k'/k$ extension \uuline{finie} \uline{séparable}, telle que $\prod_{X'/k'}P'/X'$ représentable.
\margin{\deleted{\unsure{Attention à} \unsure{l'énoncé} : \unsure{lorsqu'on} \illegible{2} $P/X$ \unsure{diff. exemples}, mais $P/X$ de type fini, l'énoncé est-il \uline{vrai} ???} \deleted{??}}

\inserted{\uline{Corollaire 4} Si $P/X$ à fibres quasi-projectives, et $X/k$ propre \deleted{\unsure{plat}}, alors $\prod_{X/k}P/X$ existe et est ess. quasi-proj. (\unsure{Ceci|Ainsi} : Cor. 1) \deleted{\unsure{de}}}

\uline{Corollaire 5} $P/X/S$, $X$ propre et plat sur $S$ noeth intègre, $P/X$ loc. de type fini. Alors $\exists$ \deleted{voisinage} ouvert $U\ne\varnothing$, et morphisme fini étale surjectif $U'\unsure{\longrightarrow|\text{—}}U$, tel que $\prod_{X_{U'}/U'}P_{U'}/X_{U'}$ représentable. Si $X_{\unsure{\eta(k(U))|\eta(U)}}$ satisfait la condition énoncée dans Cor. 2, on peut même prendre $U'=U$.
\margin{\deleted{\illegible{1} Faut-il supposer $P/X$ \unsure{diff. ?} \unsure{intègre} ??}}
% archive foliation: 26
